\section{Topology}
\label{sec:topology}

This section describes how \vannak{} places durable metadata across \ipto{}
instances. Placement is deterministic, epoch-versioned, and replicated through
\raft{}, while the ring itself is rebuilt deterministically on each node.

\subsection{Shard Model}

\vannak{} distinguishes two unrelated notions of ``shard'':

\begin{description}
\item[\sitas{} shard id] where work is currently processed (executor placement).
\item[\texttt{DataIndividualShardId}] which \ipto{} repository instance owns the
  durable metadata for a data item (domain placement).
\end{description}

A \sitas{} shard may ingest, buffer, index, or retry a metadata event, but the
final \ipto{} destination is selected by \texttt{DataIndividualShardId} through
the placement resolver. The two must never be conflated.

\subsection{Shard ID Derivation}

\texttt{DataIndividualShardId} is a \texttt{u64} newtype. A shard id can be
derived deterministically from a stable data-individual identity using
\texttt{DataIndividualShardId::from\_data\_individual}, which hashes the
identity string with an FNV-1a derivative so ids are uniformly distributed
across the \texttt{u64} space.

\begin{lstlisting}
pub fn from_data_individual(data_individual_id: &DataIndividualId) -> Self {
    let mut state = 0xcbf2_9ce4_8422_2325u64;
    for byte in data_individual_id.as_str().as_bytes() {
        state ^= u64::from(*byte);
        state = state.wrapping_mul(0x0000_0100_0000_01b3);
    }
    Self(state)
}
\end{lstlisting}

Callers may instead assign an explicit \texttt{DataIndividualShardId} when the
placement key carries business meaning, such as scoping all data for a tenant to
a dedicated shard range.

\subsection{Consistent Hash Ring}

The primary placement mechanism is a deterministic consistent-hash ring over the
\texttt{u64} space, represented by \texttt{IptoPlacementRing}. Each \ipto{}
instance contributes a configurable number of virtual nodes. The ring is sorted
and binary-searched to resolve a \texttt{DataIndividualShardId} to an instance.

\begin{lstlisting}
pub fn resolve(&self, shard_id: DataIndividualShardId) -> Option<&IptoInstanceId> {
    if self.points.is_empty() {
        return None;
    }
    let idx = self
        .points
        .binary_search_by_key(&shard_id.0, |p| p.position)
        .unwrap_or_else(|i| i);
    Some(&self.points[idx % self.points.len()].instance)
}
\end{lstlisting}

\subsection{Virtual Nodes}

Each \texttt{IptoPlacementSlot} pairs an \texttt{IptoInstanceId} with a
\texttt{vnodes} count. A slot must have at least one virtual node; constructing a
slot with zero virtual nodes returns \texttt{ClusterControlError::ZeroVnodes}.
More virtual nodes per instance produce a smoother distribution. Typical
configurations use 64 to 256 virtual nodes per instance.

\subsection{FNV-1a Hash Placement}

Virtual node positions are computed by \texttt{ring\_hash(instance, vnode\_idx)},
the same FNV-1a family used for the segment checksum and shard derivation. Both
the instance name and the virtual node index are fed through the full mixing
loop, giving good distribution and full determinism.

\begin{lstlisting}
fn ring_hash(instance: &str, vnode_idx: u32) -> u64 {
    let mut state = 0xcbf2_9ce4_8422_2325u64;
    for byte in instance.as_bytes() {
        state ^= u64::from(*byte);
        state = state.wrapping_mul(0x0000_0100_0000_01b3);
    }
    for byte in &vnode_idx.to_le_bytes() {
        state ^= u64::from(*byte);
        state = state.wrapping_mul(0x0000_0100_0000_01b3);
    }
    state
}
\end{lstlisting}

\subsection{Ring Determinism}

Because the ring is built only from the slot configuration and a fixed hash
function, every node constructs an identical ring from the same slots. The ring
points themselves are never replicated; \raft{} replicates only the slots and
overrides. Two rings built from the same slots resolve every shard id
identically.

\subsection{Adding and Removing Instances}

Adding or removing an \ipto{} instance changes only a small fraction of the
shard space, proportional to one over the number of instances. Empirically,
growing a three-instance ring to four instances (256 virtual nodes each) moves
roughly one quarter of shards, and the moved fraction is bounded well away from
a full reshuffle.

\subsection{Range Overrides}

Operators may pin an inclusive shard range to a specific instance with an
\texttt{IptoPlacementRange}. Overrides take priority over the ring, so they are
the mechanism for deliberate, business-driven placement (for example, routing a
compliance-sensitive shard range to a dedicated instance).

\begin{lstlisting}
pub struct IptoPlacementRange {
    pub start: DataIndividualShardId,
    pub end: DataIndividualShardId,
    pub target: IptoInstanceId,
}
\end{lstlisting}

A range with \texttt{end} less than \texttt{start} is rejected with
\texttt{ClusterControlError::InvalidPlacementRange}, and overlapping ranges are
rejected with \texttt{ClusterControlError::OverlappingPlacementRanges}.

\subsection{Placement Maps}

An \texttt{IptoPlacementMap} combines a ring (built from slots) with optional
range overrides, tagged by a \texttt{PlacementEpoch}. Resolution checks
overrides first, then falls back to the ring.

\begin{lstlisting}
pub fn resolve(&self, shard_id: DataIndividualShardId) -> Option<&IptoInstanceId> {
    for range in &self.overrides {
        if range.contains(shard_id) {
            return Some(&range.target);
        }
    }
    self.ring.resolve(shard_id)
}
\end{lstlisting}

A map requires at least one slot; an empty slot list is rejected with
\texttt{ClusterControlError::NoPlacementSlots}.

\subsection{Placement History and Fallback}

\texttt{ClusterControlState} keeps a bounded history of placement maps. The
constant \texttt{MAX\_PLACEMENT\_HISTORY} is 5, so the five most recent epochs
are retained and older epochs are pruned. New maps must have a strictly greater
epoch than the current one, or they are rejected with
\texttt{ClusterControlError::StalePlacementEpoch}.

The method \texttt{resolve\_with\_fallback} returns candidate instances in
priority order: the current epoch first, then previous epochs newest-first,
deduplicated.

\begin{lstlisting}
pub fn resolve_with_fallback(&self, shard_id: DataIndividualShardId)
    -> Vec<IptoInstanceId>
{
    let mut candidates = Vec::new();
    for map in self.placement_maps.values().rev() {
        if let Some(target) = map.resolve(shard_id)
            && !candidates.iter().any(|c: &IptoInstanceId| c == target)
        {
            candidates.push(target.clone());
        }
    }
    candidates
}
\end{lstlisting}

\subsection{Query Routing During Reconfiguration}

When a placement map changes, a data item may have durable metadata on its
previous owner while new writes go to the current owner. Queries use
\texttt{resolve\_with\_fallback} to try the current owner first and then prior
owners, so a reader can still find metadata written before the reconfiguration.
As a last-resort broadcast, \texttt{all\_ipto\_instances} returns every instance
known across all retained maps.

\begin{vnote}
Because history is bounded to five epochs, fallback only covers reconfigurations
within that window. Long-lived data crossing more than five placement epochs
must be rebalanced explicitly (see Section~\ref{sec:rebalancing}).
\end{vnote}

\endinput
